\documentclass[11pt]{article}
\usepackage[margin=1in]{geometry}
\usepackage{amsmath,amssymb,amsthm}
\usepackage{enumitem}
\usepackage[hidelinks]{hyperref}
\theoremstyle{plain}
\newtheorem{theorem}{Theorem}
\newtheorem{lemma}{Lemma}
\newtheorem{proposition}{Proposition}
\newtheorem{corollary}{Corollary}
\theoremstyle{remark}
\newtheorem{remark}{Remark}
\newtheorem{assumption}{Assumption}
\DeclareMathOperator{\Var}{Var}
\DeclareMathOperator{\sgn}{sgn}
\newcommand{\dFh}{\widehat{\Delta F}}
\newcommand{\dF}{\Delta F}
\newcommand{\Oh}{O(3)}
\newcommand{\SOh}{SO(3)}
\setlength{\parskip}{4pt}
\title{\vspace{-2em}The BAR Information Bottleneck\\[2pt]
\large Self-calibration (sandwich variance), $O(1)$ differentiation, a Fisher--resistance\\ correspondence, an estimation--detection conservation law, and a chirality completeness theorem}
\author{\normalsize working notes: proofs for Paper 1 (v3, corrected)}
\date{}
\begin{document}
\maketitle
\vspace{-2.5em}

\paragraph{Setup.} Two thermodynamic states $A,B$ (units $k_BT{=}1$). Forward works $\{W^f_i\}_{i=1}^{n_f}\sim P_f$, reverse works $\{W^r_j\}_{j=1}^{n_r}\sim P_r$. Unified coordinates $x_i=W^f_i$ (forward), $x_j=-W^r_j$ (reverse). By Crooks $P_f(x)/P_r(x)=e^{x-\dF}$, so the classes differ by a logistic tilt of intercept $\dF$. With $M=\ln(n_f/n_r)$, $\sigma(z)=(1+e^{-z})^{-1}$, $p(x;d)=\sigma(x-d+M)$, the Bennett estimator $\dFh$ is the MLE of the intercept (slope fixed at $1$) under the \emph{stratified two-sample} design (fixed $n_f,n_r$).

\begin{lemma}[Score, curvature, and overlap]\label{lem:fisher}
\[
S(d)=\sum_{j\in r}p(x_j)-\sum_{i\in f}(1-p(x_i)),\qquad
\partial_d S=-\!\sum_{\text{all}}p(1-p)=:-I(d).
\]
$I(d)>0$ (each $p(1-p)\in(0,\tfrac14]$), so $\ell$ is strictly concave and the root $\dFh$ is unique. $I$ is the \emph{curvature} of the objective (the ``overlap'': large iff many $p\approx\tfrac12$). It plays the role of the bread in the sandwich variance and the scale of the backward pass; it is \emph{not} in general the inverse variance (see Theorem~\ref{thm:calib}).
\end{lemma}
\begin{proof}
Stationarity is the Bennett equation; $\partial_d^2\ell=-\sum p(1-p)$ via $\sigma'=\sigma(1-\sigma)$ and $\partial_d(x-d+M)=-1$.
\end{proof}

\begin{theorem}[$O(1)$ differentiable backward; information-share gradient]\label{thm:backward}
With $\dFh(\theta)$ defined by $S(\dFh,\theta)=0$ and $\partial_d S=-I\neq0$, the implicit function theorem gives
$d\dFh/d\theta=\tfrac1I\,\partial_\theta S$. For $\theta=\mu^f$ ($x_i=\mu^f+\xi_i$) and $\theta=\mu^r$ ($x_j=-(\mu^r+\eta_j)$),
\[
\frac{d\dFh}{d\mu^f}=\frac{I^f}{I},\qquad \frac{d\dFh}{d\mu^r}=-\frac{I^r}{I},\qquad
\Big|\tfrac{d\dFh}{d\mu^f}\Big|+\Big|\tfrac{d\dFh}{d\mu^r}\Big|=1 .
\]
The backward pass is one division by the curvature $I$. (Unaffected by the variance correction below: it depends only on the deterministic second derivative.)
\end{theorem}
\begin{proof}
$\partial_{\mu^f}\!\big[-(1-p(x_i))\big]=p(1-p)$ gives $\partial_{\mu^f}S=I^f$; $\partial_{\mu^r}x_j=-1$ gives $\partial_{\mu^r}S=-I^r$. Shares sum to one since $I^f+I^r=I$.
\end{proof}
\begin{remark}[envelope variant]
If the loss is the BAR log-likelihood itself, $\partial_d\ell=0$ at the root kills the implicit term (envelope theorem); the IFT gradient is needed only when $\dFh$ feeds a downstream loss.
\end{remark}
\begin{remark}[numerical check]
On $x^f=\{0,1\}$, $x^r=\{-0.5\}$, $M=0$: $\dFh=-0.59571$, $I^f/I=0.596825$ matches the IFT gradient to machine precision; $\partial_dS=-I$ to machine precision (Wolfram Language).
\end{remark}

\begin{assumption}\label{ass:spec}
(A1) Correct specification: predicted $P_f,P_r$ equal the truth.\quad (A2) Standard M-estimation regularity.
\end{assumption}

\begin{theorem}[Self-calibration via the \emph{sandwich} variance]\label{thm:calib}
Under Assumption~\ref{ass:spec}, $\dFh$ is an M-estimator with asymptotic variance
\[
\Var(\dFh)=A^{-1}BA^{-1}=\frac{B}{I^2},\qquad
A=-\mathbb E[S']=I,\qquad
B=\Var\!\big(S(\dF)\big)=n_f\Var_f[p]+n_r\Var_r[p].
\]
The information equality $A=B$ (which would give $\Var=1/I$) holds for prospective logistic regression but \emph{fails} under the stratified two-sample design of BAR; the naive $1/I$ over-estimates the variance. The sandwich $B/I^2$ is the correct, computable-from-samples uncertainty (it coincides with the Bennett 1976 / Shirts--Chodera estimator), so calibration remains free, via the sandwich. It reduces to $1/I$ only in the information-equality (good-overlap) limit.
\end{theorem}
\begin{proof}
$\dFh$ solves $S=0$; the one-step expansion $\dFh-\dF=S(\dF)/I(\tilde d)$ and the delta method give $\Var(\dFh)=\Var(S)/I^2$ to leading order, with $\Var(S)=B$ by independence of the two samples and $\Var[1-p]=\Var[p]$.
\end{proof}
\begin{remark}[numerical check]
Correctly specified Gaussian model, $n_f=n_r=20$, $2000$ replicates: empirical $\mathrm{SD}(\dFh)=0.160$; mean sandwich se $=0.159$ (ratio $0.994$); mean naive $1/\sqrt I=0.354$ (ratio $2.21$). The sandwich is essentially exact at this $N$; the naive formula is off by $2.2\times$ in se (Wolfram Language).
\end{remark}

\begin{proposition}[Finite-sample reliability and coverage]\label{prop:finite}
Let $\widehat V=\widehat B/\widehat I^{\,2}$ be the plug-in sandwich. With probability $\ge1-\delta$,
$|\dFh-\dF|=O\!\big(\sqrt{\log(1/\delta)/N}\big)$ and $|\widehat V/V-1|=O\!\big(\sqrt{\log(1/\delta)/N}\big)$,
using Hoeffding on the bounded score/info/variance summands and the self-Lipschitz bound $|I'(d)|\le I(d)$ (since $|1-2p|\le1$). The studentized statistic obeys
\[
\Big|\,\mathbb P\big((\dFh-\dF)/\sqrt{\widehat V}\le z\big)-\Phi(z)\,\Big|\;\le\;\frac{0.56}{\sqrt B}=O(1/\sqrt N),
\]
by Berry--Esseen, using $\mathbb E|X-\mu|^3\le\Var$ for $X\in[0,1]$ so the standardized third moment is $\le 1/\sqrt{\Var S}=1/\sqrt B$.
\end{proposition}
\begin{remark}[honest scope]
This is the sampling variance under correct physics. Force-field misspecification adds a bias the frozen layer cannot self-correct (handled by a learned $\Delta$-residual head with its own epistemic $\sigma$); network error about $(\mu,\sigma)$ is epistemic (ensemble). Total Belief variance $=B/I^2+\sigma^2_{\mathrm{ens}}+\sigma^2_\Delta$.
\end{remark}

\begin{theorem}[Fisher--resistance correspondence; gauge-invariant acquisition]\label{thm:fr}
Relative measurements $y_e=(\varphi_j-\varphi_i)+\varepsilon_e$, $\varepsilon_e\sim\mathcal N(0,V_e)$ with $V_e=B_e/I_e^2$ the per-edge sandwich variance. The Gaussian posterior precision of $\varphi$ is the weighted Laplacian
\[
L=\sum_{e=(i,j)} w_e\,(e_i-e_j)(e_i-e_j)^\top,\qquad w_e=\frac1{V_e}=\frac{I_e^2}{B_e},
\]
i.e.\ the edge conductances are the \emph{inverse sandwich variances} (equal to the overlaps $I_e$ only at the information-equality limit). Then:
\begin{enumerate}[label=(\roman*),leftmargin=2.2em,itemsep=1pt]
\item $\Var(\varphi_b-\varphi_a)=(e_b-e_a)^\top L^{+}(e_b-e_a)=\Omega_{ab}$ (effective resistance).
\item A single edge has $\Omega=V_e$, recovering the BAR sandwich variance.
\item Adding a measurement along $c$ with precision $g$ updates the contrast variance in $O(1)$, $\Omega'=\Omega/(1+g\Omega)$ (Sherman--Morrison); KG equals the decision-weighted $\Omega$-reduction.
\item $\ker L=\mathrm{span}\{\text{component indicators}\}\supseteq\mathrm{span}\{\mathbf1\}$; contrasts in $\ker L$ have infinite variance and zero reduction, so $\mathrm{KG}=0$. The acquisition is automatically gauge-invariant.
\end{enumerate}
\end{theorem}
\begin{proof}
The negative log-posterior is $\tfrac12\sum_e w_e(y_e-(\varphi_j-\varphi_i))^2$ up to the gauge; its Hessian is $L$, posterior covariance $L^{+}$ on the range, giving (i)--(ii). (iii) is Sherman--Morrison on $c^\top(\cdot)^{+}c$. For (iv), all relative measurements are invariant under per-component shifts, so those directions lie in $\ker L$.
\end{proof}
\begin{remark}[numerical check]
Triangle with conductances $(1,0.5,2)$: $\Omega_{12}=0.714286$ equals the series--parallel value; $\ker L=\mathrm{span}\{\mathbf1\}$; adding edge $(1,2)$ with $g=0.7$ gives exact $\Omega'=0.476190$, matching $\Omega/(1+g\Omega)$ to machine precision (Wolfram Language).
\end{remark}

\begin{corollary}[the unified objects, corrected]
The overlap $I_e$ governs \{curvature, backward-pass scale (Thm.~\ref{thm:backward})\}. The precision $I_e^2/B_e$ governs \{calibration (Thm.~\ref{thm:calib}), Laplacian weight (Thm.~\ref{thm:fr})\}. These coincide (a single object, $V=1/I$) only at the information-equality limit. They diverge \emph{most in the high-overlap regime} ($I/B$ grows large as $p\to\tfrac12$; numerically $I/B\approx17$ at $0.5\sigma$ separation), so the naive $1/I$ over-estimates the variance most for the well-overlapped, low-variance edges, converging to the sandwich only as overlap $\to0$. The sandwich is therefore essential precisely on the \emph{reliable}, well-overlapped edges an active learner exploits. Naive $1/I$ would systematically distrust the best edges (up to ${\sim}13\times$ inflation on real FEP edges). \emph{(Direction corrected; verified three ways: controlled MC, population $I/B$, real FEP edges; see \texttt{docs/results\_figA.md}.)}
\end{corollary}

\section*{Estimation--detection conservation law}

\paragraph{Setup.} Collect the per-edge relative measurements of Theorem~\ref{thm:fr} into
$y=D\varphi+\varepsilon$, where $D\in\mathbb R^{E\times N}$ is the signed incidence matrix (row $e=(i,j)$
equals $\mathbf e_j^\top-\mathbf e_i^\top$), $\varepsilon\sim\mathcal N(0,W^{-1})$,
$W=\operatorname{diag}(w_e)$, $w_e=1/V_e$. Whitening by $W^{1/2}$ gives design $\tilde D=W^{1/2}D$; the
GLS fit has orthogonal projector (hat matrix) $H=\tilde D L^{+}\tilde D^{\top}=W^{1/2}DL^{+}D^{\top}W^{1/2}$,
with $L=D^\top W D$ the weighted Laplacian (Theorem~\ref{thm:fr}), and residual-maker $M=I_E-H$
(reusing $M$ for the residual-maker, as in the main text; the log-count ratio of the estimator setup
does not appear here). The closure statistic is $X^2=z^\top M z$ in whitened coordinates $z=W^{1/2}y$.

\begin{theorem}[Estimation--detection conservation law]\label{thm:cons}
Let the \emph{curl-leverage} be $h_e=M_{ee}$. Then
\begin{enumerate}[label=(\roman*),leftmargin=2.2em,itemsep=1pt]
\item $h_e+w_e\Omega_e=1$, with $\Omega_e$ the effective resistance across edge $e$ (Theorem~\ref{thm:fr});
\item $\sum_e h_e=\operatorname{trace}M=E-\operatorname{rank}D=E-N+c=\mathrm{dof}$, the number of
independent cycles (first Betti number; $c$ connected components);
\item under a systematic shift of size $\mu$ on edge $e$ (the mean of $y$ shifted by $\mu\,\mathbf 1_e$),
$X^2$ is noncentral $\chi^2_{\mathrm{dof}}$ with noncentrality $\lambda_e=h_e\,\mu^2/V_e$.
\end{enumerate}
Hence a bridge edge ($e$ in no cycle) has $h_e=0$ and is undetectable at every $\mu$; a node-consistent
bias $b=D\psi$ maps into $\operatorname{range}(\tilde D)=\operatorname{range}(H)$ and is annihilated by
$M$, so the test is blind to it \emph{by construction} rather than by assumption; and the smallest shift
reaching unit noncentrality is $\delta^\ast_e=\sqrt{V_e/h_e}$.
\end{theorem}
\begin{proof}
(i) $H_{ee}=w_e\,(\mathbf e_j-\mathbf e_i)^\top L^{+}(\mathbf e_j-\mathbf e_i)=w_e\Omega_e$ by
Theorem~\ref{thm:fr}(i); hence $h_e=M_{ee}=1-H_{ee}=1-w_e\Omega_e$. (ii) $H$ is the orthogonal projector
onto $\operatorname{range}(\tilde D)$, so $\operatorname{trace}H=\operatorname{rank}\tilde D
=\operatorname{rank}D=N-c$ ($W^{1/2}$ is a full-rank positive diagonal, and $\operatorname{rank}D=N-c$
for a graph with $c$ components); therefore $\operatorname{trace}M=E-(N-c)$, the cycle-space dimension.
(iii) $z\sim\mathcal N(\nu,I_E)$ since $\operatorname{Cov}(y)=W^{-1}$; under the null
$\nu\in\operatorname{range}(\tilde D)$, so $M\nu=0$ and $X^2\sim\chi^2_{\mathrm{dof}}$. Shifting the mean
of $y$ by $\mu\,\mathbf 1_e$ moves $\nu$ by $s=\mu W^{1/2}\mathbf 1_e=(\mu/\sqrt{V_e})\,\mathbf 1_e$,
giving noncentrality $\|Ms\|^2=s^\top M s=(\mu^2/V_e)\,M_{ee}=h_e\mu^2/V_e$. Finally $h_e=\|M\mathbf 1_e\|^2$
(idempotent $M$), so $h_e=0\Leftrightarrow\mathbf 1_e\in\operatorname{range}(H)$, i.e.\ $\mathbf 1_e$ lies
in the cut space $\operatorname{range}(D)$ ($e$ in no cycle); $W^{1/2}D\psi\in\operatorname{range}(\tilde D)$
gives the gradient-bias annihilation; and $\delta^\ast_e$ solves $\lambda_e=1$.
\end{proof}
\begin{remark}[numerical check]
Across the $48$ OpenFE IndustryBenchmarks2024 systems with $\ge1$ independent cycle,
$\sum_e h_e=\mathrm{dof}$ holds to a maximum deviation $1.78\times10^{-14}$, and the two evaluations of
$h_e$ (from $M=I-H$ and from $1-w_e\Omega_e$) agree to machine precision; \texttt{make figLev}.
\end{remark}

\section*{Chirality completeness}

\paragraph{Setup.} A configuration is a point cloud $\{(p_i,z_i)\}$ in $\mathbb R^3$. $\SOh$ are proper rotations ($\det R=+1$); $\Oh$ adds reflections ($\det R=-1$). The enantiomer of $M$ is $M'=\sigma M$ for a reflection $\sigma$. $M$ is chiral iff $M'$ is not a proper-rigid-motion image of $M$.

\begin{lemma}[$\Oh$-invariant readouts collapse enantiomers]\label{lem:collapse}
If $f$ is $\Oh$-invariant then $f(M')=f(\sigma M)=f(M)$. Any function of pairwise distances $\|p_i-p_j\|$ (which are $\Oh$-invariant) is constant on the pair $\{M,M'\}$; hence a distance-only network is provably enantiomer-blind.
\end{lemma}

\begin{theorem}[A parity-odd pseudoscalar separates enantiomers]\label{thm:chiral}
Let $\chi(M)=\det[\,p_2-p_1,\,p_3-p_1,\,p_4-p_1\,]$. For any linear $A$, $\chi(A\!\cdot\!M)=\det(A)\,\chi(M)$. Hence $\chi$ is $\SOh$-invariant ($\det R=+1$) and parity-odd: $\chi(M')=-\chi(M)$. For chiral $M$ (non-coplanar points) $\chi(M)\neq0$, so $\chi(M')=-\chi(M)\neq\chi(M)$ and $\chi$ distinguishes the enantiomers.
\end{theorem}
\begin{proof}
$\det[A v_1,Av_2,Av_3]=\det(A)\det[v_1,v_2,v_3]$ with $v_k=p_{k+1}-p_1$; $\det(\sigma)=-1$.
\end{proof}

\begin{corollary}[completeness: one lost bit]
For generic (chiral, full-rank) configurations, the pair $\big(\{\|p_i-p_j\|\},\ \sgn\chi\big)$ is a complete $\SOh$-invariant. Pairwise distances determine the configuration up to $\Oh$, i.e.\ up to $\{M,\sigma M\}$; $\sgn\chi$ takes opposite signs on the two and selects the enantiomer. Thus $\Oh$-invariance loses exactly the chirality bit, and a pseudoscalar restores it. Here ``complete'' means complete for the enantiomer-pair (1-bit) problem on a configuration already pinned by its full set of pairwise distances; the sign of a single triple product does not by itself characterize the chirality of an arbitrary $>4$-atom molecule, which in general requires a consistent set of pseudoscalars.
\end{corollary}

\begin{remark}[implementation]
The pseudoscalar is the parity-odd $\ell{=}0$ irrep (\texttt{0o}), realized as the triple product $v_1\!\cdot\!(v_2\times v_3)=\det[v_1,v_2,v_3]$, an odd contraction of vector features. A readout must contain a \texttt{0o} channel; a parity-even readout lies in the kernel and is blind. Note that $\|v\|$ and $v_i\!\cdot\!v_j$ are $\Oh$-invariant (even), so vector features pooled into norms/dot-products remain enantiomer-blind: the \emph{odd} triple product is required.
\end{remark}

\end{document}
